% title:          Subsets avoiding n and 2n
% area:           extremal-combinatorics
% methods:        pigeonhole, construction
% difficulty:
% difficulty-ai:  easy
% status:
% review:         none
% --- history: all optional, leave EMPTY if unknown ---
% origin:         folklore
% origin-date:
% origin-number:
% also-in:        bernoulli
% taken-from:
% references:     Part (b) is the classical 'double-free set' problem: E. T. H. Wang, On double-free sets of integers, Ars Combinatoria 28 (1989) 97-100 (earliest attribution found; cited by OEIS A050292 'Maximal cardinality of a double-free subset of {1,...,n}' - https://oeis.org/A050292 and by MathWorld 'Double-Free Set' - https://mathworld.wolfram.com/Double-FreeSet.html; a(100) = 67 by the recurrence f(n) = ceil(n/2) + f(floor(n/4))); the same problem for {1,...,3000} (no subset of 2000 elements) is in R. Honsberger, From Erdős to Kiev: Problems of Olympiad Caliber, MAA 1996, p. 177 (per MathSE 855488 - https://math.stackexchange.com/questions/855488/does-1-2-ldots-3000-contain-a-subset-of-2000-integers-with-no-member-tw; book record https://books.google.com/books?id=OIpZxK8naikC&pg=PA177); general n on MathSE 3728374 - https://math.stackexchange.com/questions/3728374/what-is-the-size-of-the-largest-subset-of-the-numbers-1-to-n-such-that-no-number; two-part form (a)+(b) with n = 100: Bernoulli Competition (EPFL) 2023, Problem 1 (checked in the club's copy of the official solutions (Bernoulli 2023-2026 solutions PDF, not online)) - https://memento.epfl.ch/event/bernoulli-competition-2023
% history:        ai
\begin{problem}
Let \( A = \{1, 2, \ldots, 100\} \) be the set of integers between 1 and 100.

(a) Let \( B \subset A \) be a subset that doesn't contain two consecutive integers. What is the maximal cardinality of \( B \)?

(b) Let \( C \subset A \) be a subset such that there is no \( n \) for which \( n \) and \( 2n \) are both in \( C \). What is the maximal cardinality of \( C \)?
\end{problem}
\begin{solution}
(a) One can easily construct a legal set \( B \) with 50 elements by setting it equal to the set of odd numbers from 1 to 100, or the set of even numbers from 1 to 100. That leaves proving that one cannot do better. In order to do that, observe that \( A \) can be partitioned into 50 pairs of an odd number and the next even number, i.e. \( A = \{1, 2\} \cup \{3, 4\} \cup \ldots \cup \{99, 100\} \). \( B \) can contain at most one element of each of these pairs, so it has at most 50 elements. Therefore, the maximum possible cardinality of \( B \) is \(50\).
\\\\
\textit{Alternate solution} One can easily construct a legal set \( B \) with 50 elements by setting it equal to the set of odd numbers from 1 to 100, or the set of even numbers from 1 to 100. That leaves proving that one cannot do better. In order to do that, first let \( x_1, \ldots, x_m \) be the elements of \( B \) in increasing order. The fact that no two elements of \( B \) are consecutive implies that \( x_{i+1} \geq x_i + 2 \) for all \( i \), so \( x_m \geq x_1 + 2(m-1) \geq 2m-1 \). However, \( x_m \leq 100 \) due to being in \( A \), so it must be the case that \( m \leq 50 \). Therefore, the maximum possible cardinality of \( B \) is \(50\).
\\\\
(b) For each nonnegative integer \( i \), let \( A_i = A \cap \{(2i+1)2^j : j \in \mathbb{Z}\} \), and observe that these are disjoint and \( A = \cup_{0 \leq i \leq 49} A_i \). The requirement on \( C \) is equivalent to saying that it never contains both \((2i+1)2^j\) and \((2i+1)2^{j+1}\), so \( |C \cap A_i| \leq \lceil |A_i|/2 \rceil \) by an argument from the solution to the previous part. We can ensure that \( |C \cap A_i| = \lceil |A_i|/2 \rceil \) for all \( i \), such as by having \( C \) be the set of every element of \( A \) that can be expressed as an odd number times a power of 4. There are 50 odd numbers in \( A \), 13 odd multiples of 4, 3 odd multiples of 16, and 1 odd multiple of 64. So, \( C \) has a cardinality of \( 50 + 13 + 3 + 1 = \boxed{67} \).
\\\\
\textit{Alternate solution:} Observe that for any \( 51 \leq m \leq 100 \), if we let \( C' = C \cup \{m\} \backslash \{m/2\} \) then \( C' \) contains at least as many elements as \( C \) and does not contain both \( n \) and \( 2n \) for any \( n \). So, we can assume that \( C \) contains \( n \) for all \( 51 \leq n \leq 100 \). In this case, it cannot contain any \( 26 \leq n \leq 50 \). Then by the same logic we can assume that \( C \) contains \(\{13, \ldots, 25\}\), which forces it not to contain any \( 7 \leq n \leq 12 \). Then we can assume that it contains \(\{4, 5, 6\}\), which means it does not contain 2 or 3, at which point we can have it contain 1. So, the maximum possible cardinality of \( C \) is \( 50 + 13 + 3 + 1 = \boxed{67} \).
\end{solution}
